% Capítulo 4 — Rotaciones
\chapter{Rotaciones}\label{cap:04}

\section{La idea}

Arrancamos por lo que dejó instalado la primera clase: un punto no es sus coordenadas, transformar es mover figuras y rotar es conservar distancias alrededor de un centro. Las cajas ``En palabras del profesor'' vienen de los resúmenes de clase de la cursada 2025; el resto es desarrollo nuestro para que el texto se sostenga solo.

\subsection{El punto y sus coordenadas no son lo mismo}\label{sec:punto-coordenadas}

{\small\vspace{1mm}\linkUnidad{01a_punto_quieto}{Simulación interactiva: Unidad 1 — Rotaciones I: el punto no se movió}\par}


Un punto geométrico existe independientemente de la manera elegida para describirlo. Sus
coordenadas son una lista de números obtenida después de elegir un origen, unos ejes, una
orientación y una unidad. Si cambia el sistema de coordenadas, cambian los números, pero no
necesariamente cambia el punto.

\begin{idea}[Idea conceptual]
La materia estudia simultáneamente dos problemas distintos:
\begin{itemize}
  \item mover o deformar objetos geométricos;
  \item cambiar la manera de asignarles coordenadas.
\end{itemize}
Las fórmulas pueden parecer similares, pero la interpretación activa o pasiva determina qué
está cambiando realmente.
\end{idea}

\paragraph{Primera traducción algebraica.}
Una posición puede representarse mediante un vector \(x\). Cambiar de descripción significa
cambiar sus coordenadas; transformar activamente el objeto significa producir un nuevo punto
\(x'=F(x)\). La fórmula es breve, pero su interpretación depende de qué permanece fijo:
el objeto, los ejes o ambos.

\begin{idea}[En palabras del profesor]
``La rotación genera para cada punto una trayectoria o línea de los lugares posibles a
los que caer. [\ldots] Estas trayectorias posibles son lo que conocemos normalmente como
líneas coordenadas'' (resumen de clase 1, cursada 2025).
\end{idea}

\subsection{Qué significa transformar una figura}\label{sec:transformar}

Una transformación asigna a cada punto una imagen. Para comprenderla geométricamente conviene
imaginar una figura sencilla, como un triángulo, y observar:
\begin{itemize}
  \item hacia dónde se desplaza cada vértice;
  \item si los lados conservan longitud;
  \item si los ángulos se mantienen;
  \item si las rectas paralelas continúan siendo paralelas;
  \item si existe un punto, una recta o una dirección que permanezca fija.
\end{itemize}

\begin{center}
\begin{tikzpicture}[scale=.88]
  \coordinate (a) at (1,1);
  \coordinate (b) at (3.4,1.2);
  \coordinate (c) at (2,3.2);
  \coordinate (ap) at (7,1.4);
  \coordinate (bp) at (10.2,1.9);
  \coordinate (cp) at (8.2,4);
  \draw[teal,line width=1.1pt] (a)--(b)--(c)--cycle;
  \draw[rust,line width=1.1pt] (ap)--(bp)--(cp)--cycle;
  \foreach \p/\q in {a/ap,b/bp,c/cp}
    \draw[auxline,-{Latex[length=2mm]}] (\p)--(\q);
  \node[labelback,text=teal] at (2.1,.55) {figura original};
  \node[labelback,text=rust] at (8.7,1.0) {figura imagen};
\end{tikzpicture}
\end{center}

La matriz y el vector de traslación son una descripción compacta de ese movimiento. Antes de
operar con ellos, debe poder anticiparse cualitativamente la figura resultante.

\paragraph{Primera traducción algebraica.}
\[
  F(x)=Ax+t.
\]
Aquí \(A\) transforma direcciones y formas, mientras que \(t\) desplaza toda la figura sin
deformarla. Las columnas de \(A\) muestran dónde terminan los vectores unitarios de los ejes.

\begin{idea}[En palabras del profesor]
``Hay transformaciones que se definen completamente por qué propiedades de los objetos
transformados mantienen'' (resumen de clase 1, rúbrica, cursada 2025).
\end{idea}

\subsection{Rotar: conservar distancias alrededor de un centro}\label{sec:rotar}

{\small\vspace{1mm}\linkUnidad{01b_rotacion_plano}{Simulación interactiva: Unidad 2 — Rotaciones II: la rotación en el plano}\par}


Una rotación mueve cada punto sobre una circunferencia cuyo centro es el centro de rotación.
Por eso conserva:
\begin{itemize}
  \item la distancia de cada punto al centro;
  \item las distancias entre puntos;
  \item los ángulos y la forma;
  \item la orientación.
\end{itemize}

En dos dimensiones basta un centro y un ángulo. En tres dimensiones también se necesita un
eje: la dirección del eje permanece fija y todo plano perpendicular a él gira como un plano
bidimensional.

\paragraph{Primera traducción algebraica.}
\[
  x'=c+R(\theta)(x-c).
\]
El vector \(x-c\) expresa la posición vista desde el centro; \(R(\theta)\) lo gira, y sumar
\(c\) devuelve el resultado al sistema original.

\begin{resultado}[Lectura inmediata]
Si una transformación 3D es una rotación, debe existir una dirección que no cambie. Esa
dirección es el eje y más adelante (sección~\ref{sec:eje-exp} y capítulo~\ref{cap:08}) aparecerá algebraicamente
como un autovector de autovalor \(1\).
\end{resultado}

\begin{idea}[En palabras del profesor]
``Una rotación implica dejar quieto un punto, mover otro, y en base a esos acomodar todos
los demás. [\ldots] Se podían acomodar los demás porque había una estructura de algo que
se conservaba (distancias y ángulos)'' (resumen de clase 1, cursada 2025).
\end{idea}

\subsection{Componer transformaciones con centros distintos}\label{sec:componer-centros}

La clase práctica combina dos operaciones:
\begin{enumerate}
  \item una rotación de \(45^\circ\) alrededor de un punto \(c\);
  \item una escala uniforme de razón \(k=0.5\) alrededor de otro punto \(d\).
\end{enumerate}
La primera cambia la orientación sin cambiar el tamaño. La segunda reduce todas las distancias
al centro \(d\) a la mitad. Como los centros son distintos, la transformación final no puede
entenderse mirando solamente una de las dos operaciones.

\paragraph{Primera traducción algebraica.}
\[
F_1(x)=c+R_{45^\circ}(x-c),
\qquad
F_2(y)=d+0.5(y-d).
\]
La composición es
\[
F(x)=F_2\bigl(F_1(x)\bigr).
\]
La expresión situada más a la derecha actúa primero: \(F_1\) rota y después \(F_2\) escala.
Al desarrollar,
\[
F(x)=0.5R_{45^\circ}x+t,
\]
donde el vector \(t\) reúne el efecto de trasladar hacia ambos centros y regresar.

\begin{idea}[Qué clase de transformación resulta]
La parte lineal \(0.5R_{45^\circ}\) gira \(45^\circ\) y reduce a la mitad. Por tanto, la
composición completa es una similaridad directa contractiva. Posee un único punto fijo, que
actúa como centro global de la transformación compuesta aunque no coincida, en general, con
ninguno de los dos centros originales.
\end{idea}

\paragraph{Movimiento discreto y trayectoria imaginada.}
Matemáticamente, \(F\) lleva cada punto directamente a su imagen. La trayectoria ``primero
un arco y después un segmento recto'' es una interpolación útil para comprender las dos
etapas, pero no forma parte de la definición de la transformación compuesta. Otra
interpolación podría unir los mismos extremos por un camino diferente.

\begin{idea}[En palabras del profesor]
``Experimentamos qué pasaba cuando aplicábamos rotaciones y traslaciones sucesivamente a
un punto (formar círculos) y tratamos de pensar por qué esto era inevitable, con un
argumento nuevamente sobre el punto fijo de la transformación'' (resumen de clase 3, cursada 2025).
\end{idea}

\section{Construcciones, cuentas y ejemplos}

Pasamos a las cuentas: la matriz de rotación en el plano y en el espacio, el número complejo como giro y el experimento de la primera clase, el que se armó en una planilla con \texttt{MMULT}. Lo dictado marca el orden y el desarrollo agrega los pasos para repetirlo con lápiz.

\subsection{Transformaciones del plano}\label{sec:transf-plano}

Una transformación del plano es una función
\[
  F:\R^2\longrightarrow\R^2.
\]
Estudiarla significa responder cuatro preguntas:
\begin{itemize}
  \item ¿qué puntos permanecen fijos?;
  \item ¿qué magnitudes conserva: distancias, ángulos, áreas o paralelismo?;
  \item ¿cómo se representa algebraicamente?;
  \item ¿qué ocurre al componerla con otra transformación?
\end{itemize}

\subsubsection{Transformaciones lineales y afines}

Una transformación lineal \(L\) cumple
\[
  L(x+y)=L(x)+L(y),\qquad L(\lambda x)=\lambda L(x).
\]
En particular, necesariamente \(L(0)=0\). Por eso una rotación alrededor del origen puede ser
lineal, pero una rotación con centro distinto del origen no puede serlo en esas coordenadas.

Una transformación afín agrega una traslación:
\[
  F(x)=Ax+t.
\]
La matriz \(A\) determina la deformación lineal y el vector \(t\) desplaza el resultado.

\subsection{Rotaciones en el plano}\label{sec:rot-plano}

{\small\vspace{1mm}\linkUnidad{01a_punto_quieto}{Simulación interactiva: Unidad 1 — Rotaciones I: el punto no se movió}\par}
{\small\vspace{1mm}\linkUnidad{01b_rotacion_plano}{Simulación interactiva: Unidad 2 — Rotaciones II: la rotación en el plano}\par}


\subsubsection{Matriz de rotación}

La rotación antihoraria de ángulo \(\theta\) alrededor del origen es
\[
  R(\theta)=
  \begin{pmatrix}
    \cos\theta&-\sin\theta\\
    \sin\theta& \cos\theta
  \end{pmatrix}.
\]
Aplicada a \(x=(x_1,x_2)^T\):
\[
  R(\theta)x=
  \begin{pmatrix}
    x_1\cos\theta-x_2\sin\theta\\
    x_1\sin\theta+x_2\cos\theta
  \end{pmatrix}.
\]

\begin{center}
\begin{tikzpicture}[scale=1.0]
  \draw[axisline,-{Latex[length=2mm]}] (-.3,0)--(5.2,0) node[right] {$x$};
  \draw[axisline,-{Latex[length=2mm]}] (0,-.3)--(0,4.4) node[above] {$y$};
  \coordinate (p) at (4,1);
  \coordinate (q) at (2.2,3.5);
  \draw[auxline] (0,0)--(p);
  \draw[auxline] (0,0)--(q);
  \draw[mainline,-{Latex[length=2.6mm]}] (p)
    arc[start angle=14,end angle=58,radius=4.12];
  \fill[teal] (p) circle (2.5pt) node[below right] {$x$};
  \fill[rust] (q) circle (2.5pt) node[above right] {$R(\theta)x$};
  \node[text=teal] at (2.1,.75) {$\theta$};
\end{tikzpicture}
\end{center}

\begin{resultado}[Propiedades]
\[
R(\theta)^TR(\theta)=I,\qquad
\det R(\theta)=1,\qquad
R(\alpha)R(\beta)=R(\alpha+\beta).
\]
La ortogonalidad implica que conserva producto interno, longitudes y ángulos. El determinante
positivo indica que también conserva la orientación.
\end{resultado}

\paragraph{Ejemplo.}
Para \(\theta=90^\circ\):
\[
R\left(\frac{\pi}{2}\right)=
\begin{pmatrix}0&-1\\1&0\end{pmatrix},
\qquad
R\left(\frac{\pi}{2}\right)(3,1)^T=(-1,3)^T.
\]

\subsubsection{Rotación alrededor de un centro cualquiera}

Para rotar alrededor de \(c\), primero se traslada el centro al origen, luego se rota y
finalmente se deshace la traslación:
\[
  \boxed{Q_{c,\theta}(x)=c+R(\theta)(x-c)}.
\]
Equivalentemente,
\[
  Q_{c,\theta}(x)=R(\theta)x+\bigl(I-R(\theta)\bigr)c.
\]
Es una transformación afín. Si \(c\neq0\), el término independiente suele ser distinto de
cero y la transformación no es lineal.

\begin{center}
\begin{tikzpicture}[scale=.95]
  \coordinate (c) at (3,2.2);
  \coordinate (p) at (6.2,1.5);
  \coordinate (q) at (5.4,4.3);
  \draw[auxline] (c)--(p);
  \draw[auxline] (c)--(q);
  \draw[mainline,-{Latex[length=2.6mm]}] (p)
    arc[start angle=-12,end angle=42,radius=3.28];
  \fill[ink] (c) circle (2.7pt) node[below left] {$c$};
  \fill[teal] (p) circle (2.5pt) node[below right] {$x$};
  \fill[rust] (q) circle (2.5pt) node[above right] {$Q_{c,\theta}(x)$};
  \draw[muted,dotted] (0,0)--(7.6,0);
  \draw[muted,dotted] (0,0)--(0,5);
  \fill[muted] (0,0) circle (1.8pt) node[below left] {$0$};
\end{tikzpicture}
\end{center}

\subsection{Rotaciones y números complejos}\label{sec:rot-complejos}

Identificando \((x,y)\in\R^2\) con \(z=x+iy\in\mathbb C\), la fórmula de Euler
\[
  e^{i\theta}=\cos\theta+i\sin\theta
\]
permite escribir una rotación como
\[
  z'=e^{i\theta}z.
\]
Si se multiplica por un complejo general \(\rho e^{i\theta}\), el módulo \(\rho\) produce
una escala uniforme y el argumento \(\theta\) produce una rotación:
\[
  z'=\rho e^{i\theta}z.
\]

\begin{idea}[Relación entre la forma compleja y la matricial]
\[
(a+ib)(x+iy)=(ax-by)+i(bx+ay)
\]
equivale a
\[
\begin{pmatrix}x'\\y'\end{pmatrix}
=
\begin{pmatrix}a&-b\\b&a\end{pmatrix}
\begin{pmatrix}x\\y\end{pmatrix}.
\]
La multiplicación compleja es una familia especial de transformaciones lineales reales:
solo usa dos parámetros, frente a los cuatro de una matriz real \(2\times2\) general.
\end{idea}

\subsection{Rotaciones en tres dimensiones}\label{sec:rot-3d}

{\small\vspace{1mm}\linkUnidad{01c_rotaciones_3d}{Simulación interactiva: Unidad 3 — Rotaciones III: tres ejes y el orden importa}\par}


\subsubsection{Rotaciones elementales}

Las rotaciones alrededor de los ejes coordenados se representan por
\[
R_x(\alpha)=
\begin{pmatrix}
1&0&0\\0&\cos\alpha&-\sin\alpha\\0&\sin\alpha&\cos\alpha
\end{pmatrix},
\]
\[
R_y(\beta)=
\begin{pmatrix}
\cos\beta&0&\sin\beta\\0&1&0\\-\sin\beta&0&\cos\beta
\end{pmatrix},
\qquad
R_z(\gamma)=
\begin{pmatrix}
\cos\gamma&-\sin\gamma&0\\
\sin\gamma& \cos\gamma&0\\
0&0&1
\end{pmatrix}.
\]

\subsubsection{El orden importa}\label{sec:orden-importa}

En general,
\[
  R_x(\alpha)R_y(\beta)R_z(\gamma)
  \neq
  R_z(\gamma)R_y(\beta)R_x(\alpha).
\]
La matriz situada más a la derecha actúa primero. Esta falta de conmutatividad fue uno de los
experimentos principales de la primera clase.

\paragraph{Ejemplo sencillo.}
Con giros de \(90^\circ\), aplicar primero \(R_x\) y luego \(R_y\) al vector
\(e_z=(0,0,1)^T\) produce
\[
R_y\left(\frac{\pi}{2}\right)
R_x\left(\frac{\pi}{2}\right)e_z
=R_y\left(\frac{\pi}{2}\right)(-e_y)
=-e_y.
\]
Al invertir el orden:
\[
R_x\left(\frac{\pi}{2}\right)
R_y\left(\frac{\pi}{2}\right)e_z
=R_x\left(\frac{\pi}{2}\right)e_x
=e_x.
\]
Como \(-e_y\neq e_x\), la diferencia no es un error numérico: es una propiedad geométrica
de las rotaciones espaciales.

\subsection{Experimento de la primera clase}

Se construyeron en Excel o LibreOffice las matrices \(R_x,R_y,R_z\) y se compararon distintos
órdenes con la función \texttt{MMULT}. Después se aplicó muchas veces una composición de
giros pequeños e iguales alrededor de los tres ejes. El resultado aproxima una rotación
alrededor de la dirección \((1,1,1)\), capaz de permutar cíclicamente los ejes:
\[
  e_x\longmapsto e_y,\qquad
  e_y\longmapsto e_z,\qquad
  e_z\longmapsto e_x.
\]
El ángulo de esa permutación es \(120^\circ\), no \(90^\circ\).

Para cuantificar el experimento, puede tomarse un paso
\[
Q_\varepsilon
=R_z(\varepsilon)R_y(\varepsilon)R_x(\varepsilon),
\qquad \varepsilon=0.01.
\]
A primer orden,
\[
Q_\varepsilon
\approx I+\varepsilon(A_x+A_y+A_z).
\]
La suma de generadores corresponde al eje \((1,1,1)\), cuya norma es \(\sqrt3\). Después de
\(n\) pasos, el ángulo acumulado es aproximadamente
\[
  \Theta\approx n\varepsilon\sqrt3.
\]
Con \(n=121\):
\[
  \Theta\approx121(0.01)\sqrt3
  \approx2.096\ {\rm rad}
  \approx120.1^\circ.
\]
La rotación exacta de \(120^\circ\) alrededor de \((1,1,1)\) que realiza la permutación
anterior tiene matriz
\[
P=
\begin{pmatrix}
0&0&1\\
1&0&0\\
0&1&0
\end{pmatrix}.
\]
Las columnas muestran directamente que \(Pe_x=e_y\), \(Pe_y=e_z\) y \(Pe_z=e_x\).

\begin{idea}[Práctica recomendada]
\begin{enumerate}
  \item Construir \(R_x,R_y,R_z\) en una hoja de cálculo.
  \item Comparar dos órdenes con ángulos grandes.
  \item Repetir con \(\varepsilon=0.01\) radianes.
  \item Calcular en Python \(\log(R)\) con \texttt{scipy.linalg.logm}.
  \item Interpretar la matriz antisimétrica obtenida.
\end{enumerate}
\end{idea}

\section{En profundidad}

Quedan dos temas que se apoyan en lo dictado y que completamos con desarrollo propio: qué pasa cuando los ángulos son chicos y por qué todo giro del espacio tiene un eje. Los apéndices~\ref{ap:infinitesimales} y~\ref{ap:rot3d} los llevan más lejos, y el eje es la puerta a los autovectores del capítulo~\ref{cap:08}.

\subsection{Rotaciones pequeñas}

Para un ángulo pequeño \(\varepsilon\),
\[
  R_x(\varepsilon)\approx I+\varepsilon A_x,
  \qquad
  R_y(\varepsilon)\approx I+\varepsilon A_y.
\]
Entonces
\[
(I+\varepsilon A_x)(I+\varepsilon A_y)
=I+\varepsilon(A_x+A_y)+\varepsilon^2A_xA_y,
\]
mientras que el orden inverso cambia solo el término de segundo orden:
\[
(I+\varepsilon A_y)(I+\varepsilon A_x)
=I+\varepsilon(A_x+A_y)+\varepsilon^2A_yA_x.
\]
Por eso las rotaciones muy pequeñas parecen conmutar a primer orden, aunque no conmuten
exactamente.

\subsection{Eje de rotación, autovectores y exponencial}\label{sec:eje-exp}

\subsubsection{El eje como dirección fija}

Toda rotación propia de \(\R^3\) posee un eje. Si \(v\) apunta en la dirección de ese eje,
\[
  Rv=v.
\]
Por lo tanto, \(v\) es un autovector de \(R\) con autovalor \(1\). Equivalentemente,
\[
  (R-I)v=0.
\]

\subsubsection{Exponencial de matrices}

La exponencial matricial se define por la serie
\[
  e^A=I+A+\frac{A^2}{2!}+\frac{A^3}{3!}+\cdots.
\]
Si una rotación se expresa como \(R=e^A\), entonces
\[
  R^t=e^{tA},\qquad 0\le t\le1,
\]
describe una trayectoria continua desde la identidad hasta \(R\). Del mismo modo,
\[
  e^{A/n}
\]
es un paso pequeño que, aplicado \(n\) veces, produce \(e^A\).

\begin{resultado}[Generador antisimétrico]
Para una rotación tridimensional, el logaritmo matricial principal produce, bajo condiciones
adecuadas, una matriz antisimétrica
\[
A=
\begin{pmatrix}
0&-\omega_z&\omega_y\\
\omega_z&0&-\omega_x\\
-\omega_y&\omega_x&0
\end{pmatrix}.
\]
El vector \(\omega=(\omega_x,\omega_y,\omega_z)\) indica el eje y su norma representa el
ángulo de rotación.
\end{resultado}
